\documentclass[conference]{IEEEtran}
\usepackage{cite}
\usepackage[spanish]{babel}
\selectlanguage{spanish}
\usepackage[utf8]{inputenc}
\hyphenation{op-tical net-works semi-conduc-tor}


\begin{document}
\title{Plantilla de Survey para CS401 \\
Aquí poner el título de tu Survey } 


% author names and affiliations
% use a multiple column layout for up to three different
% affiliations
\author{\IEEEauthorblockN{Alumno 1}
\IEEEauthorblockA{ Ciencia de la Computación\\Universidad Católica San Pablo\\
Email: ejemplo@correo.com}
\and
\IEEEauthorblockN{Autor 2}
\IEEEauthorblockA{Ciencia de la Computación\\Universidad Católica San Pablo\\
Email: ejemplo@correo.com}
}



% make the title area
\maketitle

% As a general rule, do not put math, special symbols or citations
% in the abstract
\begin{abstract}
Es importante contar con una plantilla para el desarrollo del survey elaborado en el curso de MICC, con el fin de guiar la redacción del mismo. 
En el abstract se escribe un breve resumen de la investigación que han realizado.
\end{abstract}

% no keywords

\begin{IEEEkeywords} Plantilla \LaTeX, ProyectoI \end{IEEEkeywords}

\IEEEpeerreviewmaketitle

\section{Introducción}

Colocar aquí el contexto, motivación, objetivo y estructura del $Survey$.


\section{Marco Teórico}
Colocar aquí los conceptos que serán usados dentro del $Survey$ para un mejor entendimiento.


\section{Estado del Arte}
En los primeros párrafos el alumno deberá mencionar a los \textit{surveys} relacionados con el tema.
Luego, deberá elaborar una taxonomía (clasificación) PROPIA de los trabajos revisados (dentro de esta categoría no entran \textit{surveys}) y mostrarlo en una tabla resumen y/o figura. 

Explicar cada clase de la taxonomía en una subsección diferente.

\subsection{Clase 1 de la Taxonomía- Generar un nombre}
Colocar aquí la descripción general de la clase así como la importancia de la misma. \\
Por cada uno de los artículos que pertenecen a esta clase, hacer una descripción (objetivo,metodología, validación de resultados, fortalezas y debilidades de la técnica) 


\subsection{Clase n de la Taxonomía}
Colocar aquí la descripción general de la clase así como la importancia de la misma. \\
Por cada uno de los artículos que pertenecen a esta clase, hacer una descripción (objetivo,metodología, validación de resultados, fortalezas y debilidades de la técnica)\cite{HuangEtAl2007}.

El último párrafo de esta sección hará referencia a una tabla resumen de los artículos revisados según la taxonomía propuesta.

\section{Conclusiones }
Redactar conclusiones claras y concisas a las que ha llegado después de hacer una exhaustiva investigación sobre el tema.


\bibliographystyle{IEEEtran} 
\bibliography{biblio}
% completa y estandarizada

\end{document}
